This study asked how current evidence on sleep-dependent memory consolidation,
Targeted Memory Reactivation (TMR), and wearable sleep monitoring can be
translated into a scientifically defensible first-generation pathway for
real-time closed-loop TMR. The evidence supports a qualified but actionable
answer. TMR is no longer only a biologically plausible hypothesis: controlled
experiments and meta-analytic evidence demonstrate a modest, condition-dependent
effect on subsequent memory, while home-based studies show that physiologically
guided cueing can be automated outside the conventional sleep laboratory
\citep{hu2020promoting,whitmore2022improving,salfi2025promoting}. At the same
time, the literature does not establish a mature memory-enhancement technology,
a universally effective cueing protocol, or an integrated system whose
reliability has been demonstrated across users, tasks, devices, and repeated
nights.

The principal translational conclusion is therefore the selection of an
\textbf{EEG-guided, stage-aware, real-time closed-loop TMR software layer,
intended for integration with compatible existing wearable EEG hardware, as
the first-generation reference pathway}. This selection is not based on evidence
that EEG is permanently superior to peripheral sensing. Peripheral physiology
has already supported automated home TMR, and later reduced-sensor or hybrid
approaches remain scientifically credible. EEG is selected first because it
currently preserves closer continuity with the electrophysiological domain used
to characterize sleep and TMR, while avoiding the additional inference required
by peripheral-only sensing. Stage-aware rather than phase-locked control follows
the same reasoning: the evidence for state-dependent NREM cueing is broader than
the evidence that finer oscillatory targeting produces a reproducible incremental
behavioral advantage. Integration with existing hardware similarly represents a
sequencing decision intended to avoid combining intervention validation with
unnecessary hardware-development uncertainty.

The broader contribution of this work is not the specification of a finished
device, algorithm, or product architecture, but the construction of a traceable
evidence-to-engineering argument. Biological findings were translated into
functional implications; those implications were converted into system-level
requirements; candidate implementation pathways were then compared according to
physiological observability, temporal suitability, evidence maturity, and the
additional assumptions required for validation. This process identified a
critical distinction that extends beyond TMR: the existence of individually
feasible components does not demonstrate the validity of their integration.
Retrospective sleep-stage accuracy does not establish intervention readiness,
technical execution does not establish behavioral efficacy, and engineering
convenience does not substitute for physiological justification. A defensible
neurotechnology pathway must preserve these distinctions rather than resolve
scientific uncertainty through design choice alone.

The resulting pathway therefore remains a testable reference configuration,
not a validated technology. Prospective work must establish that usable EEG can
be maintained during overnight operation, that physiological eligibility can be
estimated causally with appropriate handling of uncertainty, that cue delivery
remains temporally linked to the state that authorized it, that stimulation does
not systematically destabilize sleep, and that technically faithful operation
produces a reproducible behavioral effect. Failure at that final step would be
scientifically meaningful: if the system reliably reproduces the intended
physiological intervention yet repeated controlled studies fail to reproduce a
memory benefit, increasing algorithmic complexity would not by itself justify
continued translation. The intervention hypothesis, task conditions, or
practical significance of the pathway would require reconsideration.

Accordingly, the endpoint of this study is neither a claim of deployment
readiness nor an argument for technological finality. It is a narrowing of the
translational problem to a configuration in which the remaining uncertainties
can be observed, tested, and potentially falsified. If validated, the EEG-guided
reference pathway could provide the experimental baseline against which
adaptive control, finer temporal targeting, multimodal sensing, reduced-channel
EEG, and peripheral-only approaches are evaluated. The central outcome is
therefore a scientifically bounded route from evidence to experiment: one that
defines not only what should be investigated first, but also what evidence must
exist before the next technological claim is justified.